% !TeX encoding = UTF-8
\documentclass[variant=apuntes,cover=institutional,mode=screen]{ua-engineering-v6-article}
% !TeX encoding = UTF-8
\UASetUniversity{Universidad Austral}
\UASetFaculty{Facultad de Ingeniería}
\UASetDepartment{Departamento de Computación}
\UASetCampus{Pilar}
\UASetCareer{Ingeniería Informática}
\UASetCommission{Comisión A}
\UASetProfessor{Equipo docente}
\UASetSemester{Segundo cuatrimestre 2026}
\UASetCourse{Sistemas Distribuidos}
\UASetDate{2026}


\UASetClassNumber{11}
\UASetClassTitle{Resiliencia distribuida}
\UASetDocKind{Apuntes}

\title{Clase 11 --- Resiliencia distribuida}
\author{\UAProfessor}
\date{\UADate}

\begin{document}
\maketitle

\section{Introducción}

En sistemas distribuidos, la falla no es una excepción rara, sino una condición normal del entorno. La cuestión central no es si el sistema fallará, sino \emph{cómo} fallará y \emph{qué tan lejos} se propagará esa falla.

Una dependencia lenta, un broker saturado o una base de datos bajo presión pueden transformar una degradación local en un colapso sistémico. Por eso la resiliencia distribuida no consiste en “agregar reintentos” sin más, sino en diseñar mecanismos que contengan la propagación del daño.

\begin{quote}
Resiliencia es la capacidad del sistema de sostener comportamiento útil bajo estrés, degradación o falla parcial.
\end{quote}

\section{Fallos locales y fallos sistémicos}

Una falla local es un problema en un componente específico:
\begin{itemize}
\item un servicio responde lento;
\item una cola se llena;
\item una base de datos supera su capacidad;
\item un cache deja de responder.
\end{itemize}

El verdadero problema aparece cuando el resto del sistema reacciona mal:
\begin{itemize}
\item multiplica requests;
\item acumula trabajo sin límite;
\item consume recursos compartidos;
\item propaga latencia a otras capas.
\end{itemize}

Esto transforma una falla localizada en una falla sistémica.

\section{Retries correctos}

\subsection{Intuición}

Retry es útil cuando la falla es transitoria. Ejemplos:
\begin{itemize}
\item una pérdida de paquete;
\item jitter corto de red;
\item una condición temporal recuperable.
\end{itemize}

\subsection{Problema}

Retry es peligroso cuando la causa es:
\begin{itemize}
\item saturación;
\item cola excesiva;
\item falla persistente;
\item operación no idempotente.
\end{itemize}

En esos casos, reintentar agrega más presión sobre el componente que ya está en problemas.

\subsection{Buenas prácticas}

\begin{itemize}
\item límite estricto de intentos;
\item exponential backoff;
\item jitter para evitar sincronización patológica;
\item aplicar solo en operaciones seguras o idempotentes;
\item evitar retry en múltiples capas a la vez.
\end{itemize}

\subsection{Razón profunda}

Un retry no es gratis. Cada retry consume:
\begin{itemize}
\item ancho de banda;
\item conexiones;
\item CPU;
\item slots de concurrencia;
\item tiempo de espera en colas.
\end{itemize}

\section{Circuit Breakers}

\subsection{Idea}

Un circuit breaker corta temporalmente las llamadas a una dependencia degradada cuando detecta demasiadas fallas o latencias excesivas.

\subsection{Estados}

\paragraph{Closed}
Las requests fluyen normalmente.

\paragraph{Open}
Se rechazan llamadas sin intentar ejecutarlas.

\paragraph{Half-open}
Se deja pasar un pequeño conjunto de requests de prueba para evaluar si la dependencia se recuperó.

\subsection{Ventajas}

\begin{itemize}
\item evita insistir sobre una dependencia caída o saturada;
\item protege recursos locales;
\item reduce tiempo desperdiciado en espera inútil;
\item da feedback más rápido al sistema llamador.
\end{itemize}

\subsection{Peligros}

\begin{itemize}
\item thresholds mal ajustados;
\item abrir demasiado pronto;
\item reabrir demasiado tarde;
\item usarlo sin estrategia de fallback.
\end{itemize}

\section{Backpressure}

\subsection{Idea}

Cuando el consumidor no puede sostener el ritmo, el productor debe desacelerarse o el sistema debe rechazar trabajo.

\subsection{Sin backpressure}

\begin{itemize}
\item las colas crecen;
\item la latencia sube;
\item aumenta el uso de memoria;
\item aparecen timeouts;
\item se disparan retries;
\item el sistema entra en espiral.
\end{itemize}

\subsection{Con backpressure}

\begin{itemize}
\item el ingreso se regula;
\item se rechaza trabajo antes del colapso;
\item el sistema protege su estabilidad.
\end{itemize}

\subsection{Formas de backpressure}

\begin{itemize}
\item límites de cola;
\item ventanas de crédito;
\item control explícito de demanda;
\item rate limiting;
\item rechazo temprano.
\end{itemize}

\section{Bulkheads}

\subsection{Analogía}

En ingeniería naval, un bulkhead separa compartimentos para que una vía de agua no hunda todo el barco. En sistemas distribuidos, el patrón hace lo mismo con recursos.

\subsection{Aplicación}

\begin{itemize}
\item pools de hilos separados por dependencia;
\item colas independientes;
\item límites de concurrencia por tipo de request;
\item aislamiento por tenant o dominio.
\end{itemize}

\subsection{Valor}

Si una dependencia consume todos los hilos, conexiones o buffers del sistema, otras rutas críticas quedan bloqueadas. Bulkheads evitan esa contaminación transversal.

\section{Graceful Degradation}

\subsection{Idea}

Cuando el sistema no puede ofrecer la funcionalidad completa, intenta ofrecer una versión degradada pero todavía útil.

\subsection{Ejemplos}

\begin{itemize}
\item servir catálogo sin recomendaciones;
\item responder desde cache aun si está desactualizada;
\item desactivar features no críticas;
\item aceptar solo lecturas;
\item rechazar operaciones caras mientras se sostiene el núcleo del servicio.
\end{itemize}

\subsection{Propósito}

Graceful degradation permite que el sistema falle en forma parcial y controlada, en lugar de caer completamente.

\section{Combinación de patrones}

La resiliencia real no suele depender de un único mecanismo. Una estrategia razonable combina:

\begin{itemize}
\item timeouts bien elegidos;
\item retries limitados con backoff;
\item circuit breakers;
\item aislamiento de recursos;
\item control de presión;
\item degradación funcional;
\item observabilidad suficiente para decidir.
\end{itemize}

\section{Error común}

Un error muy habitual es introducir solo retries y timeouts, sin circuit breakers ni backpressure. Eso crea un sistema que parece más robusto, pero en realidad responde al estrés con más carga, más espera y más amplificación de la falla.

\section{Modelo del curso}

Podemos expresar el problema con:

\[
D = (P, F, G, S, T)
\]

Ejemplo:
\begin{itemize}
\item $P$: mantener servicio útil bajo degradación parcial;
\item $F$: dependencia lenta, timeouts, saturación, colas crecientes;
\item $G$: contención de la falla y preservación parcial del servicio;
\item $S$: retries limitados + circuit breaker + bulkheads + backpressure + degradación;
\item $T$: más complejidad de configuración, tuning y operación.
\end{itemize}

\section{Conclusión}

La resiliencia distribuida no elimina la falla. La vuelve:
\begin{itemize}
\item contenida;
\item observable;
\item recuperable;
\item menos destructiva.
\end{itemize}

\begin{uakeybox}
Un sistema resiliente no es el que nunca falla, sino el que impide que una falla local se transforme en colapso sistémico.
\end{uakeybox}


\section{Síntesis razonada de la clase}
Esta unidad debe leerse como parte de una cadena de decisiones de diseño y no como un catálogo de mecanismos. Para estudiar el tema con profundidad conviene reconstruir cuatro preguntas: \emph{qué problema aparece}, \emph{qué garantía se desea}, \emph{qué mecanismo la aproxima} y \emph{qué costo introduce}. En cada ejemplo debe identificarse además el modelo de falla implícito.

\begin{uakeybox}
Una respuesta técnicamente madura no termina al nombrar un algoritmo o patrón: declara sus supuestos, la propiedad que preserva y el escenario en el que deja de ser suficiente.
\end{uakeybox}

\section{Bibliografía guiada y ruta de profundización}
Para profundizar esta clase se recomienda: Google SRE Book; Nygard, Release It!; AWS Builders Library sobre retries y jitter.

La lectura debe hacerse con preguntas concretas: (1) ¿qué modelo de sistema asume el autor?, (2) ¿qué propiedad demuestra o persigue?, (3) ¿qué falla o anomalía motiva el mecanismo?, y (4) ¿qué trade-off queda abierto?

\section{Conexión con el Trabajo Final Integrador}
El grupo debe señalar qué parte de su sistema utiliza los conceptos de esta unidad, justificar por qué son necesarios y documentar una alternativa descartada. Cuando el contenido todavía no corresponda a la etapa vigente, debe registrarse como hipótesis o decisión pendiente.

\section{Autoevaluación conceptual}
\begin{enumerate}
\item ¿Cuál es el problema distribuido concreto que motiva esta clase?
\item ¿Qué propiedad de seguridad y qué propiedad de progreso están en juego?
\item ¿Qué supuesto de red, tiempo o falla resulta decisivo?
\item ¿Qué trade-off aceptaría en un sistema real y cuál no aceptaría en un dominio crítico?
\item ¿Cómo observaría experimentalmente que la solución se comporta como espera?
\end{enumerate}

\section{Profundización autocontenida v9}
\subsection{Problema, límite e invariante}
El núcleo de esta clase es retry budgets, backoff, jitter, circuit breakers, backpressure, bulkheads y graceful degradation. Para estudiarlo de manera autónoma conviene separar tres planos. Primero, el \emph{problema observable}: qué comportamiento incorrecto puede ver un cliente o un operador. Segundo, el \emph{límite del modelo}: qué información, sincronía o coordinación no está disponible. Tercero, el \emph{invariante}: qué propiedad no debe violarse aunque el sistema pierda progreso temporalmente.

El modelo de fallas relevante incluye dependencia lenta, retry storm, cola sin límite, pool agotado y cascada. Una solución debe describirse como protocolo y no solo como nombre de tecnología: actores, estados, mensajes, condición de éxito, condición de aborto y estado visible después de una interrupción. La evidencia esperada es curvas de carga, colas, breaker state y degradación del camino crítico.

\subsection{Método de análisis}
Ante un caso nuevo, siga esta secuencia:
\begin{enumerate}
\item Identifique actores, dato o recurso compartido y frontera de confianza.
\item Escriba una ejecución nominal y una ejecución adversarial.
\item Distinga seguridad (lo malo no ocurre) de progreso (algo bueno finalmente ocurre).
\item Declare qué supuesto permite avanzar: mayoría, deadline, sincronía parcial, operación idempotente, almacenamiento durable u otro.
\item Compare una alternativa más simple y una más fuerte; registre qué métrica o experimento haría cambiar la decisión.
\end{enumerate}

\subsection{Caso transferible}
Considere una plataforma de reservas que recibe operaciones duplicadas, pierde conectividad entre regiones y debe sostener un camino crítico sin ocultar el estado real al usuario. Aplique el mecanismo de esta clase a una única decisión concreta. Una respuesta madura no intenta resolver todo: delimita la operación, formula la garantía y explica la degradación esperada cuando el supuesto central deja de cumplirse.

\subsection{Errores de razonamiento frecuentes}
\begin{itemize}
\item confundir una propiedad con el producto que suele implementarla;
\item describir el camino feliz sin recorrer una falla intermedia;
\item afirmar ``exactly-once'', ``alta disponibilidad'' o ``consistencia'' sin definir alcance observable;
\item medir solo throughput aunque la hipótesis trate sobre semántica o recuperación;
\item presentar la complejidad añadida como beneficio sin compararla con la alternativa más simple.
\end{itemize}

\section{Lectura guiada}
Nygard, Release It!; Google SRE Workbook; patrones de backpressure y load shedding.

Durante la lectura, registre: modelo de sistema, propiedad perseguida, contraejemplo que motiva el mecanismo, costo principal y una condición bajo la cual no lo elegiría. Estas cinco anotaciones convierten la bibliografía en una herramienta de diseño y no en una lista para memorizar.

\section{Aplicación al Trabajo Final Integrador}
El equipo debe producir un ADR breve con la decisión asociada a esta clase. Debe contener: contexto, dos alternativas, supuesto decisivo, garantía elegida, trade-off, evidencia pendiente y fecha de revisión. La evidencia mínima esperada para esta unidad es: curvas de carga, colas, breaker state y degradación del camino crítico.

\section{Autoevaluación avanzada}
\begin{enumerate}
\item ¿Puedo explicar la clase sin nombrar primero una tecnología?
\item ¿Puedo construir una ejecución que viole la garantía si el mecanismo se configura mal?
\item ¿Puedo distinguir seguridad, progreso y experiencia del usuario?
\item ¿Puedo proponer una medición que valide la propiedad, no solo el rendimiento?
\item ¿Puedo defender cuándo una solución más simple sería suficiente?
\end{enumerate}

\end{document}
